\documentclass{article}
\usepackage{tikz}
\usepackage{xcolor}
\usepackage{amsmath}
\usepackage{fontspec}
\usepackage{xeCJK}  % 添加对中文的支持

% 设置中文字体 - 根据系统可用字体调整
\setCJKmainfont{STSong}  % 或其他可用中文字体

\usetikzlibrary{shapes,arrows,calc,positioning,fit,backgrounds}

\definecolor{convcolor}{RGB}{245,230,200}
\definecolor{bncolor}{RGB}{220,240,220}
\definecolor{relucolor}{RGB}{230,230,250}
\definecolor{attncolor}{RGB}{255,204,204}
\definecolor{inoutcolor}{RGB}{204,229,255}
\definecolor{poolcolor}{RGB}{255,255,204}
\definecolor{mulcolor}{RGB}{230,210,255}
\definecolor{blockcolor}{RGB}{240,240,255}

\begin{document}
\pagestyle{empty} % 确保没有页眉页脚，最大化内容区域
\begin{center}
\begin{tikzpicture}[
    node distance=0.6cm, % 减小默认节点间距
    box/.style={rectangle, draw, minimum width=2.7cm, minimum height=0.8cm, text centered, rounded corners=2pt, font=\footnotesize},
    convbox/.style={rectangle, draw, fill=convcolor, minimum width=2.7cm, minimum height=0.8cm, text centered, font=\footnotesize},
    bnrelubox/.style={rectangle, draw, fill=bncolor, minimum width=2.7cm, minimum height=0.6cm, text centered, font=\footnotesize},
    relubox/.style={rectangle, draw, fill=relucolor, minimum width=2.7cm, minimum height=0.6cm, text centered, font=\footnotesize},
    poolbox/.style={trapezium, draw, fill=poolcolor, minimum width=2.7cm, minimum height=0.6cm, text centered, trapezium stretches=true, font=\footnotesize},
    circlebox/.style={circle, draw, fill=mulcolor, minimum size=0.7cm, text centered, font=\footnotesize},
    inoutbox/.style={rectangle, draw, rounded corners=12pt, fill=inoutcolor, minimum width=2.7cm, minimum height=0.8cm, text centered, font=\footnotesize},
    arrow/.style={thick,->,>=stealth},
    line/.style={thick,-},
    block/.style={rectangle, draw, dashed, thick, fill=blockcolor, opacity=0.7, inner sep=7pt}, % 减小块内边距
    blocklabel/.style={rectangle, fill=white, text opacity=1, font=\bfseries\footnotesize} % 标签字体已设为footnotesize
]

% Input
\node (input) [inoutbox] {输入图像};

% LDS Branch
\node (conv1) [convbox, below=1.2cm of input] {Conv 3×3, 32}; % 减小与input的距离
\node (bn1) [bnrelubox, below=of conv1] {BN + ReLU};
\node (conv2) [convbox, below=0.7cm of bn1] {Conv 3×3, 64}; % 减小层间距
\node (bn2) [bnrelubox, below=of conv2] {BN + ReLU};
\node (conv3) [convbox, below=0.7cm of bn2] {Conv 3×3, 128}; % 减小层间距
\node (bn3) [bnrelubox, below=of conv3] {BN + ReLU};

% Channel Attention Module
\node (avgpool) [poolbox, right=4.0cm of bn3] {AdaptiveAvgPool2d(1)}; % 减小右移距离
\node (attn1) [convbox, above=of avgpool] {Conv 1×1, 32};
\node (relu) [relubox, above=of attn1] {ReLU};
\node (attn2) [convbox, above=of relu] {Conv 1×1, 128};
\node (sigmoid) [relubox, above=of attn2] {Sigmoid};

% Multiplication node
\node (multiply) [circlebox, below=1.0cm of bn3] {$\times$}; % 减小与bn3的距离

% Upsampling and Classifier
\node (upsample) [poolbox, below=1.0cm of multiply] {双线性插值 ×4}; % 明确与multiply的距离
\node (classifier) [convbox, below=1.0cm of upsample] {Conv 1×1, num\_classes}; % 减小与upsample的距离
\node (output) [inoutbox, below=0.8cm of classifier] {分割图}; % 减小与classifier的距离

% Connections
\draw [arrow] (input) -- (conv1);
\draw [arrow] (conv1) -- (bn1);
\draw [arrow] (bn1) -- (conv2);
\draw [arrow] (conv2) -- (bn2);
\draw [arrow] (bn2) -- (conv3);
\draw [arrow] (conv3) -- (bn3);

\draw [arrow] (bn3) -- (multiply); % LDS to multiply

\draw [arrow] (bn3.east) -- ++(2.0,0) |- (avgpool); % 调整连接到avgpool的路径

\draw [arrow] (avgpool) -- (attn1);
\draw [arrow] (attn1) -- (relu);
\draw [arrow] (relu) -- (attn2);
\draw [arrow] (attn2) -- (sigmoid);

\draw [arrow] (sigmoid.east) -- ++(1.0,0) |- (multiply.east);

\draw [arrow] (multiply) -- (upsample); % 确保此箭头存在且可见
\draw [arrow] (upsample) -- (classifier);
\draw [arrow] (classifier) -- (output);

% Blocks
\begin{scope}[on background layer]
\node [block, fit=(conv1) (bn1) (conv2) (bn2) (conv3) (bn3)] (lds_block) {};
\node [block, fit=(avgpool) (attn1) (relu) (attn2) (sigmoid)] (attention_block) {};
\node [block, fit=(upsample) (classifier)] (classifier_block) {};
\end{scope}

% Labels on top of blocks
\node [blocklabel] at ($(lds_block.north) + (0,0.35cm)$) {学习下采样 (LDS)}; % 减小标签向上偏移
\node [blocklabel] at ($(attention_block.north) + (0,0.35cm)$) {通道注意力模块}; % 减小标签向上偏移
\node [blocklabel] at ($(classifier_block.north) + (0,0.35cm)$) {分类器}; % 减小标签向上偏移

% Tensor shape information
\node [font=\scriptsize, align=center] at ($(input.south) + (0,-0.15cm)$) {[B, 3, H, W]}; % 减小字体和偏移
\node [font=\scriptsize, align=center, fill=white, inner sep=1pt] at ($(multiply.north) + (0,0.4cm)$) {[B, 128, H/4, W/4]}; % 调整偏移
\node [font=\scriptsize, align=center] at ($(output.south) + (0,-0.15cm)$) {[B, num\_classes, H, W]}; % 减小字体和偏移

\end{tikzpicture}
\end{center}
\end{document}
